\section{Limitations and Threats to Validity}
\label{sec:limitations}

The paper's third methodological commitment is that negative and incomplete
results are reported at the same altitude as positive ones. This section
discharges that commitment for the items that a reader should weigh against
every number above. We separate threats that could change a conclusion from
gaps that merely widen a confidence interval.

\subsection{Threats That Could Change a Conclusion}

\paragraph{Benchmark debt.} All increments are measured against a base-5 OLS
incumbent that a base-2 OLS beats by $0.00082$ at DM $t = -10.6$
(Section~\ref{sec:benchmark_debt}). Orderings within tables are unaffected;
levels and the significance of the smallest effects are not. The bottom four
rows of Table~\ref{tab:buckets} have increments of the same order as the debt
and should not be treated as established until the battery is re-run on the
base-2 backbone.

\paragraph{Target winsorization.} Stage 3 clips each series to rolling 5th and
95th percentile bounds. By construction this modifies 10\% of observations, and
on a heavy-tailed variance target the clipped mass is disproportionately the
volatility spikes that a forecast exists to anticipate. Models are trained on a
clipped target and evaluated against an unclipped one, and the transform is
lossy and therefore not inverted at evaluation
(Section~\ref{sec:winsorization}). The direction of the resulting bias is not
obvious and we have not measured it.
\pending{ablation at 5/95 versus 1/99 versus no winsorization, on the shared
panel; a referee is entitled to ask whether the dense-but-weak finding survives
when the tails are not truncated, since the strongest signal may live there}

\paragraph{The diurnal divisor is itself a forecast.} Stage 1 divides by a
\emph{rolling} 20-observation same-slot mean rather than a fixed seasonal
profile (Section~\ref{sec:diurnal}). That is a slowly adapting level estimate,
and it is multiplied back in at evaluation time. Every model in the paper
therefore inherits a common and non-trivial predictive component, which
inflates all QLIKE-versus-naive comparisons and compresses differences
\emph{between} models---meaning the small effects we report may be small partly
because the divisor has already done the easy work.
\pending{the standalone out-of-sample skill of the diurnal divisor, reported as
its own arm, so that every other number can be read net of it}

This same divisor is the leading suspect for the two anomalies in the kernel
estimate of Section~\ref{sec:kernel_caveats}---negative lobes and a
steeper-than-theory exponent---since a high-pass filter produces both. The
experiment that resolves the limitation and the experiment that tests the
structural reading are the same experiment.

\paragraph{Pairwise significance is untested.} The model-class conclusion rests
on a difference between two arms that has never been tested directly
(Section~\ref{sec:class_table}). We report a consistent sign across nine
feature sets, which is evidence, but it is not the test.

\paragraph{Multiplicity.} 98 arms are compared against one shared benchmark
with no family-wise or false-discovery adjustment, and the minimum is quoted in
two places (the model-class headline and the kNN headline). The large effects
are robust to any plausible correction; the small ones are not claimed.

\subsection{Gaps and Inconsistencies}

\paragraph{Feature accounting does not reconcile.} The paper describes 41
exogenous channels, a 359-column exogenous block, and a 529-column design
matrix. Under the stated expansion rule of $|\mathcal{L}| = 6$ rolling means
per channel, $41 \times 6 = 246 \ne 359$, and $529 - 359 = 170$ is not
accounted for by six backbone columns plus calendar.
\pending{an explicit channels~$\to$~lag-expansion~$\to$~design-columns table,
per bucket; the three counts in circulation cannot all be right}

\paragraph{Two training-window conventions.} The preprocessing and backtest
description specifies a 500-observation rolling window; the causally tuned
battery uses 24{,}000 bars. Only the latter produced the numbers reported here.
The former should be removed or explicitly scoped to the superseded
experiments.

\paragraph{Panel heterogeneity in superseded tables.} Earlier subgroup analyses
compared rows with materially different $N$ (as low as 138{,}279 against
222{,}058) and differenced them anyway. Buckets whose coverage begins late---
social sentiment and options order flow---span a different calendar era from
the baseline they were differenced against, so those increments confounded
information with regime. All numbers in this version come from the frozen
battery, on which every arm shares $n = 218{,}934$ identical rows; we note the
issue because it explains why earlier internal numbers disagree with these.

\paragraph{Structural claims are not identified.} We do not fit a Hawkes
process. Section~\ref{sec:hawkes} establishes a correspondence between our
regression and a Hawkes intensity, and reports quantities ($\hat\beta$, and the
branching-ratio identity of Equation~\eqref{eq:branching_identity}) that would
be structural under that correspondence. The correspondence is not yet
established: $\hat\beta$ is a descriptive log-log slope without a standard
error, the fitted kernel violates the non-negativity that the linear Hawkes
class requires, and the decisive preprocessing experiment has not been run. We
use ``power-law-like'' rather than structural language throughout for this
reason.

\paragraph{A rival benchmark is missing.} A fractional-kernel forecaster sits
in the same power-law family as our backbone (Section~\ref{sec:rough}) and is
absent from the battery. Given the paper's one-benchmark discipline this is the
single most defensible addition to the ladder sweep, and it is inexpensive.

\subsection{Economic Significance}
\label{sec:economic}

Every effect in this paper lives between QLIKE $0.126$ and $0.134$. The joint
exogenous information effect is $0.00627$, or 4.7\% relative to the incumbent;
the model-class effect is $0.00184$, or 1.4\%; the hyperparameter oracle
ceiling is $0.00015$, or 0.1\%. We have not translated any of these into a
quantity that a user of the forecast would act on, and we think the ordering
\emph{information} $>$ \emph{model class} $>$ \emph{hyperparameters} is the
paper's transportable content precisely because the absolute magnitudes are
small and panel-specific.

Two considerations argue against over-reading the levels. First, the
calibration result of Section~\ref{sec:calibration} shows that QLIKE rewards
attenuation on this panel, so an arm's QLIKE advantage is not straightforwardly
an advantage in forecasting realized variance. Second, all reported levels are
out-of-sample in the walk-forward sense but in-sample with respect to the
choice of preprocessing pipeline, feature dictionary and evaluation window,
none of which were pre-registered.

\subsection{Reproducibility}
\label{sec:repro}

All results in this version derive from a single frozen battery of 98 arms on
one panel of $n = 218{,}934$ bars, with per-arm metrics archived alongside the
manuscript. Superseded experiments---in particular the hyperparameter search of
Section~\ref{sec:tuning_leak} and the variable-$N$ subgroup tables---are
retained in the repository history but contribute no number to this paper.
\pending{software versions, random seeds, and a code and data availability
statement}
